\clearpage\thispagestyle{empty}
\pdfbookmark[0]{Quellen zu Pallas (10. Februar und 29. August 1811)}{pallas1811letters-sources}
\begin{center}\textbf{Quellen zu Pallas (10. Februar und 29. August 1811)}\end{center}
\noindent Diese KI-Übertragung ins heutige Deutsch folgt den beiden Briefen in \href{https://gdz.sub.uni-goettingen.de/id/PPN236007467}{Gauß’ \textit{Werken}, Bd.~VI} (Göttingen, 1874), S.~331--333. Die Briefe von Carl Friedrich Gauß erschienen unter der Überschrift \textit{Astronomische Beobachtungen und Bemerkungen, vom Hrn. Prof. Gauss in Göttingen. Aus zweyen Briefen Desselben} im \textit{Astronomischen Jahrbuch für das Jahr 1814}, herausgegeben von Johann Elert Bode (Berlin, 1811), S.~\href{https://api.digitale-sammlungen.de/iiif/image/v2/bsb10538343_00254/full/1600,/0/default.jpg}{246}, \href{https://api.digitale-sammlungen.de/iiif/image/v2/bsb10538343_00255/full/1600,/0/default.jpg}{247}, \href{https://api.digitale-sammlungen.de/iiif/image/v2/bsb10538343_00256/full/1600,/0/default.jpg}{248} und \href{https://api.digitale-sammlungen.de/iiif/image/v2/bsb10538343_00257/full/1600,/0/default.jpg}{249}. Das digitalisierte Exemplar stammt aus der Bayerischen Staatsbibliothek.
\par\smallskip\noindent Der Erstdruck gibt für die Vestaopposition den 25. Mai und für die Pariser Zeit der Ceresopposition von 1807 die Stunde 3 an. In der Werkausgabe stehen der 15. Mai und die Stunde 8. Der zweite Term der Breitenstörung der Ceres beginnt im Erstdruck mit $+14{,}17$, in der Werkausgabe mit $-14.17$. Zwei Fußnoten im Erstdruck betreffen die Übermittlung der Meridianbeobachtungen an Gauß und den späteren Abdruck der Ephemeriden im selben Band. KI-Übertragung und KI-Satz.
\clearpage
